\section{Discussion}

\subsection{What AutoHeal Achieves and What It Does Not}

\textbf{Achieves:} AutoHeal eventually remediates $\geq 97.8\%$ of
high-EPSS pairs at moderate capacity (H1) under a pre-registered
six-stage architecture with locked triage thresholds and safety
bounds. The autonomous-remediation goal is met within 14 days.

\textbf{Does not achieve:} a per-window rollback rate within the
pre-registered $5\%$ tolerance (H2). The safety hard-stop fires
$14.9\%$ of the time. The autonomous pipeline frequently
escalates to human review.

\textbf{Does not achieve:} per-window throughput superiority over
Human-in-loop (H4). The conservative triage threshold caps the
AUTO bucket size at a value comparable to Human-in-loop's
unconditional capacity.

\textbf{Not analysable:} MTTR reduction (H3). The instrumentation
bug must be corrected before this hypothesis can be evaluated.

\subsection{Why H1 Succeeds Despite H4 Failing}

The H1 and H4 outcomes are non-contradictory: H1 measures
\emph{eventual} coverage at a 14-day horizon, while H4 measures
\emph{per-window} dominance over Human-in-loop. AutoHeal and
Human-in-loop converge to similar coverage by Window 12
($\bar{c}_{12} \approx 0.999$ for both); per-window AutoHeal
sometimes leads and sometimes lags, but the eventual coverage
is comparable.

The pre-registered $\geq 80\%$ dominance threshold for H4 was
calibrated to a scenario where AutoHeal's per-window throughput
substantially exceeds Human-in-loop's. In practice, with the
conservative triage threshold and safety hard-stops, AutoHeal's
per-window throughput is comparable to Human-in-loop's. H4 was
pre-registered with too aggressive an expectation; the data
falsifies that expectation cleanly.

\subsection{The Honest Operational Implication}

For a deployment under the pre-registered AutoHeal configuration:

\textit{(i) Eventual coverage is high.} A team can expect $97.8\%$
of high-EPSS CVEs remediated within 14 days at moderate capacity
(H1 supported).

\textit{(ii) Safety hard-stops are frequent.} Operations teams
should plan for the AutoHeal pipeline halting and escalating to
human review approximately $15\%$ of the time. This is not a
failure mode --- it is the safety mechanism working as designed
--- but it implies a tighter human-in-loop integration than
``set and forget'' automation suggests.

\textit{(iii) Throughput parity, not superiority.} At the
pre-registered triage thresholds, AutoHeal is not faster than
Human-in-loop. The value proposition is consistency and
safety-bounded automation, not throughput.

\textit{(iv) Threshold tuning would change the trade-off.}
Lowering the AUTO classification threshold (e.g., $0.65$ instead
of $0.80$) would widen the AUTO bucket, increase per-window
throughput, and likely flip H4. It would also lower the
average score of AUTO-acted pairs and may worsen the rollback
rate. Pre-registered evaluation of a less conservative
threshold is the natural next step.

\subsection{Why the Safety Hard-Stop Mechanism Matters}

The H2 rejection is the substantive finding for production
deployments. AutoHeal's pre-registered safety hard-stop fires
correctly when rollback rates exceed the $10\%$ threshold. Without
that mechanism, the framework would continue acting through the
rollback storm and amplify a single faulty patch into fleet-wide
operational impact.

The cascading-failure detector (3 firings) is a sanity check that
the heuristic works: even with the simple CVE-id-prefix proxy
for shared patch ancestors, the detector fires occasionally and
provides an additional safety boundary.

\textbf{Practical contribution.} AutoHeal's value as a research
artifact is not the cell-mean coverage number --- it is the
\emph{pre-registered safety mechanism that the operational data
shows firing correctly}. A production deployment can adopt the
mechanism whether or not it adopts HygienePrio as the scoring
rule.

\subsection{Connection to Papers 4--9's Findings}

The synthesis-level claims:

\textit{Paper~7's K=200 hazard (lag-1 calibration harms at high
capacity) is reproduced.} AutoHeal at $K = 200$ achieves the
same coverage as $K = 100$ but with the same halt rate and
without the calibration step (Paper~7's recipe says hold weights
fixed at $K = 200$, and AutoHeal does so).

\textit{Paper~9 Theorem~1 (closed-loop signal exhaustion) predicts
Fixed-policy K=200's drop in acted pairs from full capacity to
$4.2$/window.} The high-EPSS backlog is drained within early
windows; subsequent windows have few high-EPSS pairs to act on.
The Fixed-policy column at K=200 in Table~\ref{tab:summary}
quantifies the theorem empirically.

\textit{Paper~4~\S\ref{sec:external_validity} (synthetic-vs-real
distributional shift) does not change AutoHeal's relative ordering.}
AutoHeal's coverage trajectory under real CVE attributes
matches the synthetic-trajectory shape, with the absolute coverage
levels shifted as predicted by Paper~4.
